\documentclass[10pt,a4paper]{amsart}
\usepackage[margin=19mm]{geometry}
\usepackage{amsmath,amssymb,mathtools}
\usepackage[scr=rsfs,cal=euler]{mathalfa}
\usepackage{xfrac}
\usepackage[unicode,hidelinks,bookmarks=false]{hyperref}
\newcommand{\ie}{i.e.\ }
\newcommand{\cf}{cf.\ }
\newcommand{\resp}{respectively\ }
\newcommand{\Subsection}{\subsection}
\providecommand{\nobreakdash}{\nobreak\hbox{-}}
\setlength{\emergencystretch}{2em}
\multlinegap0pt
\usepackage{bm}
\usepackage{bbm}
\usepackage{dsfont}
\usepackage{xspace}
\usepackage{cancel}
\usepackage{mathtools}
\usepackage{amsthm}
\makeatletter
\@ifundefined{case}{\newtheorem{case}{Case}}{}
\@ifundefined{lem}{\newtheorem{lem}{Lemma}}{}
\makeatother
\title[Near-perfect matchings in highly connected 1-planar graphs w]{Near-perfect matchings in highly connected 1-planar graphs with a local crossing constraint}
\author{Licheng Zhang Yuanqiu Huang Zhangdong Ouyang}
\date{Source: arXiv:2602.05267v1 (CC BY 4.0)}
\begin{document}
\maketitle

\noindent\emph{Source and licence.} Near-perfect matchings in highly connected 1-planar graphs with a local crossing constraint. Version arXiv:2602.05267v1 (CC BY 4.0), distributed under the Creative Commons Attribution 4.0 International licence, as recorded in the arXiv licence metadata for this exact version and shown on the abstract page https://arxiv.org/abs/2602.05267v1. Full rights notice: licenses/arxiv-2602.05267-CC-BY-4.0.txt.\par\smallskip

\noindent\textit{Editorial scope.} Counted unit: the Lem whose source label is \texttt{lem:type3\_preserves\_connected} in the article (source0.tex lines 712--714, source label \texttt{lem:type3\_preserves\_connected}), delivered together with the article's own complete original proof of it (source0.tex lines 716--736, one \texttt{proof} environment of 21 lines). No mathematics is written, repaired or re-derived: statement and proof are the article's own text line for line. Required context retained uncounted: lem:localcrossing (lines 676--678). Every definition, display and label of those lines is retained unchanged. Omitted from this package and not redistributed: 1--675, 679--711, 715--715, 737--1513. Nothing inside a retained excerpt is omitted or altered: every retained body is delivered exactly as the article's lines are written, so no omission receipt of any kind accompanies this package. Citations: the retained excerpts carry no citation command at all, so no bibliography entry of the article is transcribed and no reference is redistributed. Reference adaptation: none was needed and none was made; every reference inside the delivered unit resolves inside it, so no unresolved reference or citation is produced. Printed slips kept literal: no printed word, symbol, hypothesis or number of the counted statement or of its proof was altered. Layout and class: the article's own class is \texttt{article} with packages including titlesec, amsmath, amsthm, amssymb, color, enumerate, hyperref, graphicx, tikz, bookmark, amsrefs, doi, pifont, amsgen; the wrapper is the lane's pinned \texttt{amsart} editorial wrapper at 10pt on A4, which restates the theorem-like environments the retained text needs and adds the article's own macro definitions that the retained text needs (none), with extra wrapper packages (bm, bbm, dsfont, xspace, cancel, mathtools). The delivered numbering is the unit's own. Assets: no retained span contains a figure environment, a graphics-inclusion command, a PDF-page inclusion, a table or a TikZ picture; no image file is copied into this package, the package declares no asset dependency and no image file is redistributed with it. Licence: version arXiv:2602.05267v1 of this article is distributed under the Creative Commons Attribution 4.0 International licence, as recorded in the arXiv licence metadata for this exact version and shown on the abstract page https://arxiv.org/abs/2602.05267v1. A full rights notice is delivered at licenses/arxiv-2602.05267-CC-BY-4.0.txt. Characters: every retained excerpt is pure ASCII, so the delivered engine, XeLaTeX with its default Latin Modern text font, reports no missing character.
\begin{lem}\label{lem:localcrossing}
Let $G$ be a graph with a type-2A 1-planar drawing $D$. If $x_1x_3$ crosses $x_2x_4$ in $D$, then for any $i$ where $1\le i\le 4$, $x_i$ is adjacent to $x_{i-1}$ or $x_{i+1}$ where $x_0=x_4$ and $x_5=x_1$.
\end{lem}

\begin{lem}\label{lem:type3_preserves_connected}
Let $G$ be a type-2A 1-plane graph. Let $G'$ be the graph obtained from $G$ by removing one edge from each pair of crossing edges. If $G$ is connected and nice, then $G'$ is a connected, spanning and planar subgraph of $G$.
\end{lem}

\begin{proof}
 Clearly,  $G'$ is spanning and  planar.  Suppose that $G'$ is not connected.  Since $G$ is connected, there  exist two distinct components $F^1$ and $F^2$ of $G'$ such that  $|E_G(V(F^1),V(F^2))| \ge 1$. Without loss of generality, assume that $ab \in E_G(V(F^1),V(F^2))$, where $a \in V(F^1)$ and $b \in V(F^2)$. Clearly, the edge $a b$ is crossed in $G$. We assume that $a b$ crosses with $c d$. Clearly, $c d$ is an edge of $G^{\prime}$, and thus it is impossible for $c$ and $d$ to be such that one is in $F^1$ and the other in $F^2$. Thus, we only need to consider two cases.



\begin{case} 
$c \in V(F^1)\cup V(F^2) \cup V(F)$  and $d \in V(F)$, where $F$ is a component of $G^{\prime}$ distinct from $F^1$ and $F^2$.
\end{case} 

For $c \in V(F^1)$, since $G$ is type-2A, by Lemma \ref{lem:localcrossing},  either $b c \in G^{\prime}$ or $b d \in G^{\prime}$.  The former ($b c \in G'$) contradicts that $F^1$ and $F^2$ are  distinct components, while the latter ($b d \in G'$) contradicts that $F$ and $F^2$ are  distinct components.  
Similarly, for  $c \in V(F^2)$ or $c\in V(F)$, a contradiction also arises.
 
 

\begin{case} 
 $c,d \in V(F^1)$ or  $c,d \in V(F^2)$ .  
\end{case} 
For $c,d \in V(F^1)$, since $G$ is type-2A,  by Lemma \ref{lem:localcrossing},  either $b c \in E(G')$ or $b d \in E(G')$, which contradicts that $F^1$ and $F^2$ are  distinct components.   Symmetrically, for $c,d \in V(F^2)$, a contradiction also arises. 

As the above two cases lead to contradictions, it follows that $G'$ is connected, thereby establishing the proposition.
\end{proof}

\end{document}
